The development of RoME demonstrates a successful paradigm shift in how organizations can leverage artificial intelligence for meeting intelligence. By moving beyond isolated, single-modality transcription and adopting a fused audio-visual extraction pipeline, the system achieves a far richer contextual understanding of collaborative sessions. 

The implementation of the Temporal Knowledge Graph stands as a central contribution, proving that tracking entities and decisions across time provides significantly more value than isolated meeting summaries. It transforms ephemeral conversation into a persistent, interrogatable corporate memory. Furthermore, the integration of the Model Context Protocol equips the system with true agentic autonomy, securely bridging the gap between passive summarization and active workflow execution across popular tools like Google Workspace and Slack.

As remote collaboration continues to evolve, systems like RoME will be essential for reducing administrative overhead, ensuring context is not lost between syncs, and enabling human teams to focus on high-level decision making rather than manual documentation. Future work will focus on expanding the variety of supported MCP tools and optimizing the on-device inference speed of the visual analyzer modules to support full real-time edge processing.
